\documentclass[]{article}

\usepackage{amsmath}
\usepackage{multirow}
\usepackage{natbib}
\usepackage{graphicx} 


\begin{document}

This paper investigated the conditions under which the theoretical asymptotic relationship between Eulerian spectral roughness $\xi$ and the Lagrangian fractional diffusion exponent $\alpha$, given by $\alpha = 2/\xi$, is physically observable in turbulent systems. To address this, we analyzed a hierarchy of numerical models, including multifractal energy cascades, the synthetic Kraichnan model, and the deterministic Lorenz-96 system. Our methodology combined standard Eulerian structure function analysis to determine the predicted roughness $\xi$ with a detailed Lagrangian analysis that tracked the Renormalization Group (RG) flow of an effective, time-dependent exponent $\alpha(\tau)$.

Our Eulerian analysis confirmed that the spectral roughness of the velocity fields behaved as expected, with intermittency systematically increasing the value of $\xi$. This provided a clear theoretical prediction for the corresponding Lagrangian transport. However, the analysis of tracer trajectories in the three-dimensional Kraichnan model revealed a significant departure from this prediction. Despite the theoretical expectation of strong superdiffusion, the transport remained robustly near-diffusive, with the Mean Squared Displacement scaling characteristic of Brownian motion ($H \approx 0.5$) and the displacement distributions lacking the heavy tails of a Lévy process. The RG flow analysis provided a direct visualization of this discrepancy, showing that the effective exponent $\alpha(\tau)$ remained pinned near the Gaussian value of $\alpha=2$ for all accessible time scales, failing to flow towards its predicted anomalous fixed point. This demonstrates that the system was trapped in a pre-asymptotic state, where the finite spectral resolution and simulation duration prevented the emergence of the true fractional dynamics.

Furthermore, our investigation into one-dimensional turbulent flows revealed a fundamental breakdown of the theoretical mapping. In this context, the transport was dominated by a strong, non-universal subdiffusion. This behavior is a direct consequence of topological trapping, a physical mechanism wherein particles are confined for long durations, which is not accounted for in the theoretical framework. This finding underscores that the validity of the theory is contingent on the spatial dimensionality of the system.

From these results, we have learned that while the fractional operator defined by the Eulerian roughness $\xi$ provides a valid, universal description of the asymptotic fixed point of turbulent transport, its physical manifestation is not guaranteed. The emergence of this fractional regime is critically gated by system-specific, pre-asymptotic constraints. First, a sufficient separation of scales and a long enough evolution time are required for the system to cross over from its initial near-Gaussian behavior to the anomalous state. Second, the spatial dimensionality must be sufficient to allow particles to freely sample the velocity field's fluctuations, as topological constraints in lower dimensions can introduce dominant, non-universal transport mechanisms that invalidate the theoretical link. In conclusion, the practical observability of the asymptotic fractional dynamics predicted by the Eulerian field statistics depends critically on whether the physical system can overcome these pre-asymptotic and topological barriers.

\end{document}
                